%!TEX root = ../thesis.tex
%*******************************************************************************
%****************************** Fifth Chapter **********************************
%*******************************************************************************
\chapter[Advocacy and Action]{Advocacy and Action}

% **************************** Define Graphics Path **************************
\ifpdf
    \graphicspath{{Chapter5/Figs/Raster/}{Chapter5/Figs/PDF/}{Chapter5/Figs/}}
\else
    \graphicspath{{Chapter5/Figs/Vector/}{Chapter5/Figs/}}
\fi

\epigraph{If you do not know how to ask the right question, you discover nothing.}{--- \textup{W. Edwards Deming}}


\section{Introduction}

% Talk mode about Accessible, Usable, Useful
This chapter focuses on understanding the challenges of making data \textit{useful} for the community and methods that would help guide people through the whole process of effective use of data in advocacy and action. This includes empowering people to better articulate their questions for data and helping them to come up with strategies to implement use of data.

Developing and deploying tools that promote better \textit{access}, \textit{comprehension} and \textit{usability} of data, does not mean that the data will find its way into actionable results. Data itself is not evident and inherently useful to a person or group who are able to access it. Both case studies, SenseMyStreet and Data:In Place, the technologies built aimed to make data more accessible by anchoring with place and local context for people to make sense of it. With SenseMyStreet it was the case of democratising data production and with Data:In Place it was the case of democratising data access. Both studies made good strives towards helping people access and explore different datasets on their own terms. However it might have also created an overload of choice for people, without helping people to come up with questions, queries and strategies for data use. From the case studies it became evident that it is not only up to engineers to build better tools and interfaces for data, there are also other, more human factors, that need to be explored to combine a more holistic approach for effective usage of data in civic action.

Extending on the previous case studies, this chapter discusses what makes people go from ``I want to see all the data'' to ``I want to see data that relates to a question I have'', hence providing data an actionable purpose. Building from the previous chapter, this chapter explores how we might start linking data on the web with questions that people have. After an unsuccessful deployment of added features to Data:In Place platform to address the problem, a more human centred approach was taken to facilitate the use of these features. This was done through the process of running Data Consultations in focus groups with key stakeholders from a non-profit organisation and through leveraging Right Question Institute's (RQI)\footnote{https://rightquestion.org/} Question Formulation Technique (QFT) in a community workshop setting. This chapter presents empirical evidence about the effectiveness of these methods to define actionable steps towards purposing data and enables to start mapping out the supported processes needed for effective use of data in advocacy and action. Processes that have often been neglected by designers building systems for data discovery. Additionally there is still a role of a data professional or data scientist to play in enabling the use of data by communities.

\section{Making Formal Data Request}

With the Freedom of Information Act, a new mechanism was set up for people to access data held by public authorities. Enabling everyone who wanted, to request data from public authorities by submitting a Freedom of Information (FOI) request. FOI can be looked as a formal query for data, where a person or an organisation deposits a question to the data provider (i.e., public authority). Submitting an FOI request can be a lengthy and complicated process often requiring knowledge regarding legal processes. In order to streamline these processes MySociety came up with a online platform called What They Know\footnote{\url{https://www.whatdotheyknow.com}}, that helps people to submit their requests through a web form. This has made the process of submission much easier, however the actual requesting of data has not improved. In order to get the data needed, one needs to know about the existence of some unpublished document or dataset that is being requested. Apart form using the platform to explore previous FOI requests, there is no way one can make sure that particular data is not yet published somewhere else. As long as there is no unified `Web of Data` that describes all the existing data on the internet and makes it universally accessible to any service or client, this remains an issue. Previous chapter touched on technical and organisational challenges of making the `Semantic Web` a reality.

Another scenario which can be encountered, is when a request is formulated as a direct question without an indication to particular requested resource. In that case somebody has to make the association between the question and a particular dataset (i.e., describe the relationship). Now we have reached to the same problem that was described in the previous chapter: \textit{the lack of links between questions and abstract datasets that might bare the answer}. The ways that these connections are made are largely unknown and often require specialist skills with combination of tacit knowledge. This means either understanding the organisational structures of a public authority or knowing the initial purpose for the data generation.

To unpack these issues an investigation was taken up to understand the technical approaches and social and cognitive processes linked to effective uses of data. This evolved: (1) a redesign of Data:In Place (see Chapter \ref{}) to integrate additional features for enabling people make better request for datasets that could help them answer questions or evidence local issues, and for posting challenges to get support from the larger community; (2) coming up and trailing of methods that assist people in creating questions and challenges for using data in civic action. 

\section{Dual Approach}

The aim of the case study was to explore methods and technical solutions for leveraging data in community contexts and for purposes of civic advocacy and action. The case study builds on previous work conducted in the scope of this thesis around community data production using SenseMyStreet commissioning platform and exploring the uses of open governmental data for local decision making using Data:In Place platform, and combines the knowledge gained from designing, developing, deploying and evaluation these data systems. To address the issues risen from previous case studies a more holistic approach was taken to work together with communities and implement set of \textit{mutually supporting strategies} (both technical and social cognitive) in the scope of and across two components. These include \textit{Component 1} -- environmental and systems approaches to promote usage of data and build support for data discovery and \textit{Component 2} -- interventions to make data actionable and to increase community capacity to take action. This approach simultaneously focuses on technical implementations that can be used by everyone and social cognitive strategies that focus on specific communities or publics.

This is commonly known as \textit{Dual Approach} and is widely adopted in the context of public health interventions [REFs] and increasing health equity. The Dual Approach uses a model where strategies are applied to improve the health of general population while the same time using targeted interventions to address issues encountered by a specific targeted population. It follows the principle of \textit{targeted universalism}, which states that alleviating disparities of marganilized populations is essential for developing contextually relevant strategies for achieving universal goals and improving the well-being of everyone [REF]. 

In the context of effective use of data and community action a dual approach was taken to improve the ways how people make requests for data and use strategies to help local communities deal with specific isolated questions. Table \ref{table:dual_approach} illustrates the usage of dual approach strategies applied concurrently for improving general approaches to working with and using data, and focusing on the specific targeted communities to make data actionable.

\begin{table}[ht!]
\scriptsize
\captionsetup{font=footnotesize}
\caption[Dual Approach]{Strategies for Dual Approach}% title of Table
\centering % used for centering table
\begin{tabular}{|l|c|c|c|}
\hline %insert horizontal line
\textbf{Component} & \textbf{Strategy} & \textbf{General Intervention} & \textbf{Targeted Intervention} \\[0.5ex] % inserts table heading, puts a 0.5 vertical whitespace under the text
\hline % inserts single horizontal line
% inserting body of the table
\rowcolor[gray]{.9}\textbf{Enviormental and System Approaches}	& \makecell{Build support for scalable accessing \\ and discovery of datasets} & Anchoring data requests & Data Consultations, QFT \\
\textbf{Capacity building/Building community links} & \makecell{Improve effective use \\ and actionability of data} & Systematic challenge creation & Challenge creation workshops \\ [0.5ex]
\hline%inserts single line
\end{tabular}
\label{table:dual_approach} % is used to refer this table in the text
\scriptsize
\footnotesize\emph{Source:} Table created by the author
\end{table}

% This conceptual model displays how general population strategies and interventions, and targeted interventions for priority populations, contribute to the Dual Approach.
% Working on general tools and then focusing on the specifics of the community.

% The primary goal is to reduce health disparities and improve health equity in the targeted areas.
% The primary goal was to reduce efforts of data requests and improve local community's capacity to make effective use of data for civic action.

%Component 1, environmental and lifestyle change strategies, is put into action in the same communities and jurisdictions as 

% Component 2, health systems and community-clinical linkage strategies. Local improvements are supported by state-level efforts funded by DP14-1422 and by DP13-1305 State Public Health Actions. The strategies in both components should be mutually reinforcing. 


% THE DUAL APPROACH USES STRATEGIES TO ADDRESS HEALTH ISSUES IN THE GENERAL POPULATION AT A STATE LEVEL AND THE PRIORITY POPULATIONS WITHIN TARGETED COMMUNITIES

General interventions are linked to technical features integrated into the Data:In Place and targeted interventions are methods to help local communities to understand their data needs and plan actions for data. Next two sections describe these interventions in more detail and illustrate how they could be beneficial to make  

\section{Redesign of Data:In Place}

The initial prototype of Data:In Place was build to enable community groups and charities to access and make sense of data related to their local area. The design and development of the platform was largely driven by the questions and issues of a specific public, and the datasets that were first integrated into the system were a direct response to issues and questions posited by the local Neighbourhood Planning Group at their meetings. Furthermore, additional datasets were added through engagement with local charities and looking into ways their work can be targeted and evidenced. In later stage of design of the platform a \textit{Data Request Feature} was added to the system that enabled people to put in requests for datasets digitally without having to meet face-to-face with a researcher making it a more generic for everyone to request adding datasets. Although this solved an issue physicality, it did not make it easier for people to compile these requests or for the researcher to reply to them. All these issues added to the core issue of these open datasets being underutilised by communities. Data:In place platform was facing the same issues that were pointed out in the case of What Do They Know system.

\subsection{Anchoring Data Requests}

Data:In Place case study illustrated that to make data accessible and usable to people it needs to be contextualised through place, issues and situated activities people engage in those places. These can be seen as anchors for data that make abstract datasets more comprehensible for people. However using anchoring can also be beneficial when requesting for new datasets to be added. Providing that contextual information to better understand the context where data needs to fit.

Platform's \textit{Data Request} form provided people text entries for describing the issue they are exploring, linking the actual dataset or public authority who might posses the data and describing the issue in with a story. This works well if a person has an idea about what type of data they are looking for and where to access it. Then it is a case of connecting the database with the map based querying of Data:In Place. If the particular data provider uses standardised vocabularies and descriptions of the data (such as json-schema) exposed through an API, these processes can be further atomised.

\subsection{Data Challenges}

To unpack the challenges around connecting peoples questions to abstract datasets

As we identified 
Data:In Place platform 

\section{Targeted Engagements}

\subsection{Hypothesis Setting}

With every scientific research project, first things that are posited are hypotheses that the project is trying to prove or disprove. This also entails having a clearly defined agenda and good knowledge of what kind of data and variables are needed. Case studies described in two previous chapters, although both working with different communities and datasets, had a recurring criteria for successful use of data in action. Data found use in civic action when people knew ahead or planned out what they are looking for in the data. Simply put when they were asking the right questions from the data. The scientific hypothesis approach may work in professional communities, with people who are familiar with these concepts. With the angle of democratisation these practices and this particular context however the scientific approach would not cut it. To make the process more inclusive and not serving the values of some stakeholders over others, there was a need to look at alternative approaches. But what is close to hypotheses are questions, which people are asking every day. The difference is that these are unknown and highly depended on the way they are formulated.

\subsection{Question Formulation Technique}


\section{Design of Challenge Creation Workshop}

\subsection{QFI in Community Action}

\subsection{Challenge Creation}

\section{Data Collection}

\subsection{Data Consultations}

\subsection{Community Action Day}

\section{Findings}

\section{Discussion}

\section{Conclusions}